% !PRINTED 760
\weberpage{760}

Damit ist bewiesen:

\gesperrt{Die Zahl $\pi$ ist eine transcendente Zahl. Die \glqq Quadratur des Kreises\grqq{} kann nicht durch geometrische Construction, bei der nur algebraische Curven und Flächen angewandt werden, gelöst werden.}

\begin{center}
\S.\,207.

\gesperrt{Der allgemeine Satz von Lindemann über die Exponentialfunction.}
\end{center}

Die Transcendenz der Zahlen $e$ und $\pi$, die hierdurch bewiesen ist, ist als specieller Fall in einem sehr allgemeinen Theoreme über die Exponentialfunction enthalten, das \gesperrt{Lindemann} in der oben erwähnten Abhandlung angekündigt hat, von dem ein ausgeführter Beweis in der citirten Abhandlung von \gesperrt{Weierstrass} enthalten ist. Wir wollen diesen Satz hier zum Schluss noch beweisen mit Anwendung derselben Hülfsmittel, die wir für die beiden speciellen Fälle benutzt haben. Der Satz lautet so:

I. \gesperrt{Es besteht keine Gleichung von der Form} \weq{1}{C_1e^{z_1} + C_2e^{z_2} + \cdots + C_me^{z_m} = 0,} \gesperrt{in der die Coëfficienten $C_1,C_2,\ldots,C_m$ algebraische Zahlen und die Exponenten $z_1,z_2,\ldots,z_m$ von einander verschiedene algebraische Zahlen sind, es sei denn, dass alle Coëfficienten $C_1,C_2,\ldots,C_m$ gleich Null sind.}

Um ihn zu beweisen, leiten wir zunächst einen Hülfssatz ab:

Es seien
\[
\alpha_1,\alpha_2,\ldots,\alpha_r
\]
beliebige reelle oder imaginäre, jedoch \gesperrt{von einander verschiedene} Grössen, und
\[
A_1,A_2,\ldots,A_r
\]
ebenfalls beliebige Grössen, die \gesperrt{nicht alle verschwinden}.

Dasselbe soll von den zwei Grössenreihen
\[
\begin{gathered}
\beta_1,\beta_2,\ldots,\beta_s \\
B_1,B_2,\ldots,B_s
\end{gathered}
\]
gelten. Wir bezeichnen mit
\[
\gamma_1,\gamma_2,\ldots,\gamma_t
\]
%CONT>>
